%!TeX TS-program = lualatex
%!TeX encoding = UTF-8 Unicode
%!TeX spellcheck = en-UK
% -*- coding: UTF-8; -*-
% vim: set fenc=utf-8
%=== file: myfigure.tex ==========================================
\documentclass[tikz,border=0pt]{standalone}   % “border” adds optional margin
% -------------------------------------------------------------
% 1. Packages you used in the original doc (add any that are missing)
\usepackage{amsmath,amssymb}          % math, symbols
\usepackage{tikz}
\usepackage[paperwidth=841mm,paperheight=1189mm,margin=0mm]{geometry} % A0 page size
% -------------------------------------------------------------
\begin{document}
% ------------------- START OF YOUR PICTURE --------------------
\begin{tikzpicture}[x=1mm,y=1mm]   % work in millimetres (easy for A0)

  % ---------------------------------------------------------
  % Settings
  % ---------------------------------------------------------
  \def\dotradius{0.5pt}          % radius of each dot (adjust if you want larger dots)
  \def\numdots{2000}             % how many random dots you would like
  \pgfmathsetseed{12345}        % fixed seed → reproducible picture (change or remove for true randomness)

  % ---------------------------------------------------------
  % Draw the random dots
  % ---------------------------------------------------------
  \foreach \i in {1,...,\numdots} {
    % generate a random x‑coordinate between 0 and the page width (841 mm)
    \pgfmathrandominteger{\randx}{0}{841}
    % generate a random y‑coordinate between 0 and the page height (1189 mm)
    \pgfmathrandominteger{\randy}{0}{1189}
    % place a filled circle (dot) at the random point
    \fill[black] (\randx,\randy) circle[radius=\dotradius];
  }

\end{tikzpicture}%
% -------------------- END OF YOUR PICTURE -----------------
\end{document}
%================================================================